\documentclass[11pt,a4paper]{article}

% ------------------------------------------------
% Paquetes
% ------------------------------------------------
\usepackage[utf8]{inputenc}
\usepackage[T1]{fontenc}
\usepackage[spanish,es-nodecimaldot]{babel}

\usepackage{amsmath}
\usepackage{amssymb}
\usepackage{array}
\usepackage{booktabs}
\usepackage{geometry}
\usepackage{enumitem}
\usepackage{listings}
\usepackage{xcolor}
\usepackage{float}

% ------------------------------------------------
% Configuración de página
% ------------------------------------------------
\geometry{
    margin=2.5cm
}

% ------------------------------------------------
% Configuración de código
% ------------------------------------------------
\lstset{
    basicstyle=\ttfamily\small,
    columns=fullflexible,
    keepspaces=true,
    frame=single,
    breaklines=true
}

% ------------------------------------------------
% Documento
% ------------------------------------------------
\begin{document}

\section{Especificación y procesamiento del lenguaje}

Un compilador es un programa que recibe como entrada un programa escrito en un
\textbf{lenguaje fuente} y produce como salida un programa equivalente escrito
en un \textbf{lenguaje destino}. Durante este proceso, el compilador analiza la
estructura del programa fuente, determina su significado y construye una
representación adecuada para generar el programa destino.

De acuerdo con la organización clásica presentada en el capítulo 1 del
\textit{Dragon Book}, el proceso de compilación puede dividirse en varias
fases:

\begin{enumerate}
    \item análisis léxico;
    \item análisis sintáctico;
    \item análisis semántico;
    \item generación de código intermedio;
    \item optimización de código;
    \item generación de código destino.
\end{enumerate}

Estas fases no deben entenderse necesariamente como programas independientes.
Una implementación concreta puede agrupar varias de ellas en una misma
pasada sobre el programa.

De forma general, el proceso puede representarse como:

\[
\text{Programa fuente}
\rightarrow
\text{Análisis}
\rightarrow
\text{Representación intermedia}
\rightarrow
\text{Optimización}
\rightarrow
\text{Código destino}
\]

En nuestro caso, el lenguaje destino será C, por lo que el proceso completo
terminará produciendo un programa escrito en C que posteriormente podrá ser
compilado por un compilador de C.

Como ejemplo se utilizará el siguiente programa:

\begin{lstlisting}
int x = 10 + 20;
\end{lstlisting}

\subsection{Análisis léxico}

El análisis léxico constituye la primera fase del compilador. Su función es
leer los caracteres del programa fuente y agruparlos en unidades denominadas
\textbf{lexemas}. Cada lexema es una secuencia de caracteres que pertenece a
alguna categoría léxica del lenguaje.

El analizador léxico reconoce cada lexema y produce un \textbf{token} que
representa su categoría. De esta manera, el análisis léxico establece una
separación entre la representación concreta escrita por el programador y la
categoría que interesa a las fases posteriores.

Por ejemplo, para:

\begin{lstlisting}
int x = 10 + 20;
\end{lstlisting}

podemos identificar los siguientes lexemas:

\begin{lstlisting}
int
x
=
10
+
20
;
\end{lstlisting}

Cada uno de estos lexemas pertenece a una categoría léxica:

\begin{lstlisting}
INT
IDENTIFIER
ASSIGN
INTEGER
PLUS
INTEGER
SEMICOLON
\end{lstlisting}

Por lo tanto, el análisis léxico responde principalmente a la pregunta:

\begin{quote}
\textit{¿Qué unidades léxicas aparecen en la secuencia de caracteres del
programa fuente?}
\end{quote}

El analizador léxico no determina si esas unidades pueden combinarse
correctamente. Esa responsabilidad corresponde al analizador sintáctico.

\subsection{Lexemas, tokens y patrones}

Es importante distinguir tres conceptos relacionados.

Un \textbf{lexema} es una secuencia concreta de caracteres que aparece en el
programa fuente. Por ejemplo, en:

\begin{lstlisting}
int contador = 25;
\end{lstlisting}

los lexemas incluyen:

\begin{lstlisting}
int
contador
=
25
;
\end{lstlisting}

Un \textbf{token} es la categoría léxica asignada a un lexema. Por ejemplo:

\begin{lstlisting}
int       -> INT
contador  -> IDENTIFIER
=         -> ASSIGN
25        -> INTEGER
;         -> SEMICOLON
\end{lstlisting}

Finalmente, un \textbf{patrón} es una regla que describe qué secuencias de
caracteres pueden formar lexemas pertenecientes a un determinado token.

Para nuestro lenguaje se puede definir inicialmente:

\begin{table}[H]
\centering
\begin{tabular}{>{\ttfamily}c c c}
\toprule
\textbf{Token} & \textbf{Patrón o lexema} & \textbf{Ejemplo} \\
\midrule
INT &
\texttt{int} &
\texttt{int} \\

IDENTIFIER &
\texttt{[a-zA-Z\_][a-zA-Z0-9\_]*} &
\texttt{x}, \texttt{contador} \\

INTEGER &
\texttt{[0-9]+} &
\texttt{10}, \texttt{250} \\

ASSIGN &
\texttt{=} &
\texttt{=} \\

PLUS &
\texttt{+} &
\texttt{+} \\

SEMICOLON &
\texttt{;} &
\texttt{;} \\
\bottomrule
\end{tabular}
\caption{Especificación léxica inicial del lenguaje}
\label{tab:tokens}
\end{table}

Por lo tanto, para el programa:

\begin{lstlisting}
int x = 10 + 20;
\end{lstlisting}

el analizador léxico podría producir una secuencia como:

\begin{lstlisting}
INT
IDENTIFIER(x)
ASSIGN
INTEGER(10)
PLUS
INTEGER(20)
SEMICOLON
\end{lstlisting}

Los tokens que poseen un valor asociado, como los identificadores y los
literales enteros, pueden transportar información adicional. Por ejemplo:

\begin{lstlisting}
IDENTIFIER(x)
INTEGER(10)
\end{lstlisting}

El identificador también puede representarse mediante una referencia a una
entrada de la tabla de símbolos:

\begin{lstlisting}
(id, 1)
\end{lstlisting}

donde $id$ representa la categoría del token y $1$ identifica la entrada
correspondiente en la tabla de símbolos.

\subsection{Análisis sintáctico}

El análisis sintáctico recibe la secuencia de tokens producida por el
analizador léxico y determina si dicha secuencia puede derivarse de acuerdo
con la gramática del lenguaje.

Mientras que el análisis léxico determina qué unidades aparecen en el
programa, el análisis sintáctico determina cómo pueden combinarse esas
unidades.

Por ejemplo, la secuencia:

\begin{lstlisting}
INT IDENTIFIER ASSIGN INTEGER PLUS INTEGER SEMICOLON
\end{lstlisting}

puede corresponder a una declaración válida:

\begin{lstlisting}
int x = 10 + 20;
\end{lstlisting}

En cambio:

\begin{lstlisting}
INT IDENTIFIER ASSIGN PLUS INTEGER SEMICOLON
\end{lstlisting}

contiene tokens válidos individualmente, pero su organización no corresponde
a una construcción permitida por la gramática.

Por lo tanto, el análisis sintáctico responde a la pregunta:

\begin{quote}
\textit{¿Cómo se combinan los tokens para formar construcciones válidas del
lenguaje?}
\end{quote}

\subsection{Especificación sintáctica}

La sintaxis del lenguaje se especifica mediante una gramática libre de
contexto. Para el lenguaje propuesto podemos definir inicialmente:

\begin{align*}
\langle programa \rangle
    &\rightarrow \langle declaracion \rangle \\[4pt]
\langle declaracion \rangle
    &\rightarrow
    \langle tipo \rangle
    \ \texttt{IDENTIFIER}
    \ \texttt{ASSIGN}
    \ \langle expresion \rangle
    \ \texttt{SEMICOLON} \\[4pt]
\langle tipo \rangle
    &\rightarrow \texttt{INT} \\[4pt]
\langle expresion \rangle
    &\rightarrow
    \langle expresion \rangle
    \ \texttt{PLUS}
    \ \langle termino \rangle \\[4pt]
\langle expresion \rangle
    &\rightarrow \langle termino \rangle \\[4pt]
\langle termino \rangle
    &\rightarrow \texttt{INTEGER}
\end{align*}

Esta gramática establece qué secuencias de tokens forman declaraciones y
expresiones válidas.

Para:

\begin{lstlisting}
int x = 10 + 20;
\end{lstlisting}

el analizador sintáctico puede construir una estructura equivalente a:

\begin{lstlisting}
declaracion
|-- tipo -> INT
|-- IDENTIFIER -> x
|-- ASSIGN -> =
|-- expresion
|   |-- expresion
|   |   `-- termino
|   |       `-- INTEGER -> 10
|   |-- PLUS
|   `-- termino
|       `-- INTEGER -> 20
`-- SEMICOLON
\end{lstlisting}

Esta estructura constituye una representación de la organización sintáctica
del programa.

\subsection{Árbol de sintaxis abstracta}

A partir del análisis sintáctico puede construirse un
\textbf{árbol de sintaxis abstracta} (AST). A diferencia del árbol de
derivación completo, el AST elimina detalles sintácticos que no son
necesarios para las fases posteriores y conserva principalmente la estructura
semánticamente relevante.

Para:

\begin{lstlisting}
int x = 10 + 20;
\end{lstlisting}

un AST puede representarse como:

\begin{lstlisting}
Declaration
|-- type: Int
|-- name: x
`-- value: BinaryExpression
    |-- left: IntegerLiteral(10)
    |-- operator: +
    `-- right: IntegerLiteral(20)
\end{lstlisting}

El AST permite que las siguientes fases trabajen sobre una representación
estructurada del programa en lugar de procesar directamente los caracteres
del código fuente.

Además, la estructura del árbol permite representar propiedades sintácticas
como la precedencia de operadores. Por ejemplo:

\begin{lstlisting}
10 + 20 * 3
\end{lstlisting}

debe interpretarse como:

\[
10 + (20 \times 3)
\]

y no como:

\[
(10 + 20) \times 3
\]

El AST representa esta diferencia mediante la posición de los nodos dentro
del árbol.

\subsection{Análisis semántico}

El hecho de que un programa sea sintácticamente correcto no implica que sea
válido desde el punto de vista semántico.

El análisis semántico utiliza el árbol sintáctico y la información asociada
a los identificadores para comprobar restricciones que no pueden expresarse
adecuadamente mediante la gramática.

Por ejemplo:

\begin{lstlisting}
int x = "hola";
\end{lstlisting}

puede tener una estructura sintáctica válida, ya que presenta la forma general:

\begin{lstlisting}
tipo + identificador + asignacion + expresion
\end{lstlisting}

Sin embargo, la asignación puede ser semánticamente incorrecta porque una
expresión de tipo cadena no puede asignarse directamente a una variable de
tipo \texttt{int}.

Para una declaración:

\begin{lstlisting}
int x = e;
\end{lstlisting}

podemos establecer inicialmente las siguientes reglas semánticas:

\begin{enumerate}
    \item $x$ debe ser un identificador válido.
    \item $x$ no debe haber sido declarado previamente en el mismo ámbito.
    \item La expresión $e$ debe producir un valor de tipo \texttt{int}.
    \item El valor producido por $e$ se asigna a $x$.
    \item La variable $x$ queda registrada como una variable de tipo
    \texttt{int}.
\end{enumerate}

Para la expresión:

\begin{lstlisting}
10 + 20
\end{lstlisting}

se puede determinar:

\[
10 : \texttt{int}
\]

\[
20 : \texttt{int}
\]

\[
10 + 20 : \texttt{int}
\]

Por lo tanto:

\begin{lstlisting}
int x = 10 + 20;
\end{lstlisting}

es semánticamente válida.

\subsection{Tabla de símbolos}

La \textbf{tabla de símbolos} es una estructura de datos utilizada por
distintas fases del compilador para almacenar información sobre los nombres
que aparecen en el programa.

No constituye necesariamente una fase independiente del compilador, sino un
servicio utilizado principalmente durante el análisis semántico y,
posteriormente, durante la generación de código.

Entre los atributos que puede contener una entrada se encuentran:

\begin{itemize}
    \item nombre;
    \item tipo;
    \item ámbito;
    \item ubicación;
    \item información adicional necesaria para la generación de código.
\end{itemize}

Para:

\begin{lstlisting}
int x = 10 + 20;
\end{lstlisting}

podemos obtener:

\begin{table}[H]
\centering
\begin{tabular}{ccc}
\toprule
\textbf{Nombre} & \textbf{Tipo} & \textbf{Ámbito} \\
\midrule
\texttt{x} & \texttt{int} & local \\
\bottomrule
\end{tabular}
\caption{Ejemplo de tabla de símbolos}
\label{tab:symbol-table}
\end{table}

Cuando el analizador léxico encuentra el lexema \texttt{x}, puede producir
un token asociado a una entrada de esta tabla:

\begin{lstlisting}
IDENTIFIER -> (id, 1)
\end{lstlisting}

Las fases posteriores pueden utilizar esta referencia para obtener
información sobre el identificador.

\subsection{Manejo de errores}

El compilador también debe detectar y comunicar errores encontrados durante
el procesamiento del programa.

Los errores pueden aparecer en distintas fases. Por ejemplo:

\begin{itemize}
    \item un carácter que no pertenece al lenguaje puede producir un error
    léxico;
    \item una secuencia de tokens que no puede derivarse mediante la
    gramática puede producir un error sintáctico;
    \item una operación incompatible con los tipos puede producir un error
    semántico.
\end{itemize}

El manejo de errores es una actividad transversal al proceso de compilación.
El objetivo no es solamente detectar un error, sino también proporcionar
información suficiente para identificar su ubicación y naturaleza.

\subsection{Generación de código intermedio}

Después del análisis, el compilador puede construir una
\textbf{representación intermedia} (IR).

La representación intermedia proporciona una forma de separar el lenguaje
fuente del lenguaje o arquitectura destino. En lugar de generar directamente
código específico de una máquina, el compilador representa las operaciones
del programa mediante una forma intermedia.

Para:

\begin{lstlisting}
int x = 10 + 20;
\end{lstlisting}

una representación intermedia sencilla podría ser:

\begin{lstlisting}
t1 = 10 + 20
x = t1
\end{lstlisting}

La IR representa las operaciones necesarias sin depender todavía de una
arquitectura específica.

Una ventaja importante de esta separación es que permite modificar
independientemente el front-end y el back-end del compilador.

\subsection{Optimización de código}

Una vez construida la representación intermedia, pueden aplicarse
transformaciones que preserven el significado del programa pero mejoren sus
propiedades.

Por ejemplo:

\begin{lstlisting}
t1 = 10 + 20
x = t1
\end{lstlisting}

puede transformarse en:

\begin{lstlisting}
x = 30
\end{lstlisting}

Esta transformación corresponde a la evaluación de una expresión constante,
conocida como \textit{constant folding}.

Las optimizaciones pueden realizarse sobre diferentes aspectos del programa,
como expresiones, instrucciones, flujo de control o utilización de valores.

La optimización no modifica el significado observable del programa. Su
objetivo es obtener una representación equivalente que pueda ser ejecutada
de manera más eficiente.

\subsection{Generación de código destino}

La última fase principal del compilador transforma la representación
intermedia en el lenguaje destino.

En un compilador tradicional, el lenguaje destino puede ser directamente un
lenguaje ensamblador o código máquina de una arquitectura determinada.

Conceptualmente, esta fase debe resolver cuestiones como:

\begin{itemize}
    \item qué instrucciones representan cada operación;
    \item dónde almacenar los valores;
    \item qué registros utilizar;
    \item cómo representar las operaciones de control;
    \item cómo organizar la memoria necesaria.
\end{itemize}

El proceso puede representarse como:

\[
\text{IR}
\rightarrow
\text{selección de instrucciones}
\rightarrow
\text{asignación de registros}
\rightarrow
\text{código destino}
\]

\subsection{Front-end y back-end}

Las fases del compilador pueden agruparse conceptualmente en
\textbf{front-end} y \textbf{back-end}.

El front-end depende principalmente del lenguaje fuente. Sus
responsabilidades incluyen:

\begin{itemize}
    \item análisis léxico;
    \item análisis sintáctico;
    \item análisis semántico;
    \item generación de una representación intermedia.
\end{itemize}

El back-end recibe la representación intermedia y se ocupa de transformarla
en código del lenguaje destino, incluyendo las optimizaciones y decisiones
relacionadas con la generación de código.

De forma simplificada:

\[
\text{Fuente}
\rightarrow
\boxed{\text{Front-end}}
\rightarrow
\text{IR}
\rightarrow
\boxed{\text{Back-end}}
\rightarrow
\text{Destino}
\]

Esta separación permite que diferentes lenguajes fuente compartan un mismo
back-end:

\[
\begin{array}{ccc}
\text{Lenguaje A} & \rightarrow & \text{Front-end A} \\
\text{Lenguaje B} & \rightarrow & \text{Front-end B}
\end{array}
\rightarrow
\text{IR}
\rightarrow
\text{Back-end}
\]

De forma similar, un mismo front-end podría utilizarse con distintos
back-ends:

\[
\text{Fuente}
\rightarrow
\text{Front-end}
\rightarrow
\text{IR}
\rightarrow
\begin{cases}
\text{Back-end x86} \\
\text{Back-end ARM}
\end{cases}
\]

Esta separación constituye una de las principales ventajas de utilizar una
representación intermedia.

\subsection{Fases y pasadas}

Las fases representan las distintas tareas lógicas que debe realizar el
compilador. Sin embargo, estas fases no determinan necesariamente cómo se
organiza la implementación.

Una implementación puede realizar varias fases durante una misma
\textbf{pasada} sobre el programa.

Por ejemplo, una implementación puede realizar:

\[
\text{Fuente}
\rightarrow
\text{análisis léxico}
\rightarrow
\text{análisis sintáctico}
\rightarrow
\text{análisis semántico}
\rightarrow
\text{IR}
\]

como una primera pasada.

Posteriormente puede realizar:

\[
\text{IR}
\rightarrow
\text{IR optimizada}
\]

y finalmente:

\[
\text{IR optimizada}
\rightarrow
\text{código destino}
\]

Por lo tanto, las \textbf{fases} describen las tareas conceptuales del
compilador, mientras que las \textbf{pasadas} describen cómo se organizan
esas tareas durante su implementación.

\subsection{Herramientas para la construcción de compiladores}

Las diferentes fases de un compilador pueden implementarse manualmente o
mediante herramientas especializadas.

Entre las herramientas utilizadas habitualmente se encuentran:

\begin{itemize}
    \item \textbf{generadores de analizadores léxicos}, que producen
    analizadores a partir de especificaciones de tokens y patrones;

    \item \textbf{generadores de analizadores sintácticos}, que producen
    analizadores a partir de una gramática;

    \item \textbf{motores de traducción orientados a la sintaxis}, que
    permiten asociar acciones o reglas de traducción a estructuras
    sintácticas;

    \item \textbf{generadores de código}, que facilitan la construcción de
    back-ends;

    \item \textbf{herramientas de análisis de flujo}, utilizadas para
    obtener información necesaria para distintas optimizaciones;

    \item \textbf{kits de construcción de compiladores}, que proporcionan
    componentes para implementar varias de estas tareas.
\end{itemize}

Estas herramientas no constituyen fases adicionales. Son mecanismos que
facilitan la implementación de las fases existentes.

\subsection{Proceso completo}

El proceso completo puede resumirse de la siguiente manera:

\[
\boxed{
\text{Fuente}
\rightarrow
\text{Análisis léxico}
\rightarrow
\text{Tokens}
}
\]

\[
\boxed{
\text{Tokens}
\rightarrow
\text{Análisis sintáctico}
\rightarrow
\text{AST}
}
\]

\[
\boxed{
\text{AST}
\rightarrow
\text{Análisis semántico}
\rightarrow
\text{AST validado}
}
\]

\[
\boxed{
\text{AST validado}
\rightarrow
\text{Código intermedio}
\rightarrow
\text{Optimización}
}
\]

\[
\boxed{
\text{IR optimizada}
\rightarrow
\text{Generación de código}
\rightarrow
\text{Código destino}
}
\]

En cada una de estas fases intervienen estructuras auxiliares, como la tabla
de símbolos y los mecanismos de manejo de errores.

\section{Uso de C como lenguaje destino}

En nuestro compilador se utilizará C como lenguaje destino. Esto significa
que nuestro compilador no generará directamente código máquina ni
ensamblador. En cambio, generará un programa escrito en C que represente las
operaciones del lenguaje fuente.

El proceso será:

\[
\text{Lenguaje fuente}
\rightarrow
\text{Análisis léxico}
\rightarrow
\text{Análisis sintáctico}
\rightarrow
\text{Análisis semántico}
\rightarrow
\text{IR}
\rightarrow
\text{Código C}
\rightarrow
\text{Compilador de C}
\rightarrow
\text{Código máquina}
\]

Por lo tanto, el objetivo de nuestro back-end será traducir la
representación intermedia a código C válido.

Por ejemplo, si nuestra representación intermedia contiene:

\begin{lstlisting}
t1 = 10 + 20
x = t1
\end{lstlisting}

el back-end podría generar:

\begin{lstlisting}
int x;
x = 10 + 20;
\end{lstlisting}

El compilador de C se encargará posteriormente de realizar las etapas
necesarias para producir código específico de la plataforma.

Esta decisión simplifica considerablemente la construcción del back-end.
Los compiladores de C existentes proporcionan optimizaciones maduras y
soporte para múltiples arquitecturas y plataformas. Por lo tanto, nuestro
proyecto no necesita implementar desde cero las instrucciones, registros,
convenciones y características específicas de cada arquitectura.

Si, en cambio, generáramos código máquina directamente, deberíamos
implementar un back-end específico para cada arquitectura que quisiéramos
soportar.

La utilización de C como lenguaje destino permite entonces concentrar el
trabajo del proyecto en el diseño y procesamiento de nuestro propio
lenguaje, delegando al compilador de C la generación final de código para
la plataforma correspondiente.

\section{Conclusión}

A partir de este modelo se puede definir el desarrollo de nuestro compilador
como una secuencia de etapas claramente diferenciadas.

En primer lugar, se debe definir la especificación léxica del lenguaje,
determinando sus tokens, lexemas y patrones. Luego se debe establecer la
gramática que determina cómo pueden combinarse esos tokens. Sobre esta
estructura se incorporarán las reglas semánticas y la tabla de símbolos
necesarias para validar los programas.

Una vez definido el front-end, se establecerá una representación intermedia
que permita expresar las operaciones del lenguaje de manera independiente
del código C generado. Finalmente, se implementará el back-end encargado de
traducir dicha representación a C.

De esta manera, el desarrollo seguirá la organización clásica de un
compilador:

\[
\boxed{
\text{Lenguaje}
\rightarrow
\text{Léxico}
\rightarrow
\text{Sintaxis}
\rightarrow
\text{Semántica}
\rightarrow
\text{IR}
\rightarrow
\text{C}
}
\]

Esta estructura permite construir el compilador progresivamente y mantener
separadas las responsabilidades de cada etapa.

\end{document}
